\section{Results}
\label{sec:results}

\subsection{Full-channel subject-independent performance}
With all 16 electrodes, XGBoost reached a Macro-F1 of \FullXGB\ (95\% CI \FullXGBLCB--\FullXGBUCB) over the 43 held-out participants. The RBF-SVM reached \FullSVM\ (\FullSVMLCB--\FullSVMUCB) and the random forest \FullRF\ (\FullRFLCB--\FullRFUCB) (Table~\ref{tab:full_channel}, Fig.~\ref{fig:full}). XGBoost and the SVM did not differ: the paired difference was $\XGBvsSVMDiff$ (95\% CI $\XGBvsSVMLCB$ to $\XGBvsSVMUCB$), $p_{\mathrm{Holm}} \XGBvsSVMP$, $r_{rb} = \XGBvsSVMR$. Both exceeded the random forest ($p_{\mathrm{Holm}} \XGBvsRFP$ and $\SVMvsRFP$; $r_{rb} = \XGBvsRFR$ and $\SVMvsRFR$). The SVM result is a lower bound because of its training subsample (Methods). The spread between participants was far larger than any difference between models: XGBoost ranged from \FullXGBMin\ to \FullXGBMax\ across participants.

\begin{table*}[t]
\centering
\small
\caption{Full 16-channel leave-one-subject-out performance on GRABMyo (43 participants, 16 active gestures; chance Macro-F1 $\approx$ 0.063). CI: participant cluster bootstrap.}
\label{tab:full_channel}
\begin{tabular}{llcccc}
\toprule
Model & Training windows & Macro-F1 & 95\% CI & SD & Range \\
\midrule
XGBoost & all ($\approx$550k) & 0.542 & 0.517--0.567 & 0.085 & 0.385--0.702 \\
RBF-SVM & 40k & 0.538 & 0.513--0.562 & 0.085 & 0.383--0.718 \\
Random forest & 200k & 0.502 & 0.475--0.528 & 0.090 & 0.346--0.670 \\
\bottomrule
\end{tabular}
\par\smallskip\parbox{0.95\linewidth}{\footnotesize Paired comparisons (Wilcoxon signed-rank, Holm over 3): Random forest $-$ RBF-SVM = $-$0.0362 (95\% CI $-$0.0444 to $-$0.0284), $p_{\mathrm{Holm}}$ $<$ 0.001, $r_{rb}$ = -0.97; Random forest $-$ XGBoost = $-$0.0408 (95\% CI $-$0.0475 to $-$0.0340), $p_{\mathrm{Holm}}$ $<$ 0.001, $r_{rb}$ = -0.98; RBF-SVM $-$ XGBoost = $-$0.0045 (95\% CI $-$0.0107 to +0.0015), $p_{\mathrm{Holm}}$ = 0.154, $r_{rb}$ = -0.25.}
\end{table*}


\begin{figure*}[t]
\centering
\includegraphics[width=\textwidth]{figures/fig3_full_channel.pdf}
\caption{Full-channel LOSO performance (43 participants; each point is one held-out participant; grey lines join the same participant; black bar = median; diamond = mean with 95\% cluster-bootstrap CI). (a) Three classifiers. (b) Paired XGBoost $-$ SVM difference per participant, sorted; grey band = 95\% CI of the mean difference.}
\label{fig:full}
\end{figure*}

\subsection{H1: a 12-electrode subset is non-inferior under the prespecified margin}
Channel attribution was strongly concentrated and highly reproducible across outer folds. The same four electrodes (\TopFourSet) formed the top four in \TopFourFolds\ of 43 folds, and they carried \TopFourShare\% of the summed attribution. The eight lowest-ranked electrodes carried \BottomEightShare\% (Fig.~\ref{fig:attribution}). Electrodes \DroppedFourSet\ were the four excluded at $k = 12$ in \DroppedFourFolds\ of 43 folds.

Performance declined smoothly as electrodes were removed (Fig.~\ref{fig:curve}, Table~\ref{tab:sparse_ni}). The $k = 16$ reference of the curve (\CurveKSixteen) differs from the full-data result above because every curve model uses the same fixed 200,000-window training subsample. At $\delta = 0.02$, the smallest non-inferior subset was $k^* = \KStar$ electrodes: $\bar\Delta_{12} = \DeltaKTwelve$ (95\% CI $\DeltaKTwelveLCB$ to $\DeltaKTwelveUCB$). At $k = 11$ the lower bound ($\DeltaKElevenLCB$) crossed the margin. Every $k$ from 12 to 15 was non-inferior, so the result is not an isolated crossing. For the same reason, the result survives the multiplicity of testing 15 values of $k$. A fixed-sequence procedure, testing $k = 15, 14, \ldots$ in order and stopping at the first failure, controls the familywise error rate without adjustment and yields the same $k^*$. The sensitivity margins gave $k^* = \KStarMarginPtOne$ at $\delta = 0.01$ and $k^* = \KStarMarginPtThree$ at $\delta = 0.03$. Removing four electrodes (\KStarReductionPct\% of the array) is therefore within the prespecified tolerance. Larger reductions are not: halving the array to $k = 8$ cost $\DeltaKEight$ (95\% CI $\DeltaKEightLCB$ to $\DeltaKEightUCB$), and $k = 4$ cost $\DeltaKFour$. All paired differences at $k < 16$ were negative. The claim is non-inferiority within a stated engineering tolerance, not equivalence.

\begin{table*}[t]
\centering
\small
\caption{Sparse sensor curve and non-inferiority (cross-fitted ranking, $\lambda = 0$; 43 participants; every $k$ retrained from scratch on a fixed 200k-window outer-training subsample). $\Delta_k$ = paired difference to the $k = 16$ model of the same curve. Non-inferior when the lower 95\% bound exceeds $-\delta$; $\delta = 0.02$ is primary.}
\label{tab:sparse_ni}
\begin{tabular}{rcccccc}
\toprule
$k$ & Macro-F1 (mean $\pm$ SD) & $\Delta_k$ & 95\% CI & $\delta$=0.01 & $\delta$=0.02 & $\delta$=0.03 \\
\midrule
16 & 0.541 $\pm$ 0.084 & reference &  &  &  &  \\
15 & 0.538 $\pm$ 0.083 & $-$0.0036 & $-$0.0056 to $-$0.0017 & yes & yes & yes \\
14 & 0.535 $\pm$ 0.083 & $-$0.0060 & $-$0.0081 to $-$0.0039 & yes & yes & yes \\
13 & 0.530 $\pm$ 0.083 & $-$0.0109 & $-$0.0132 to $-$0.0085 & no & yes & yes \\
12 & 0.528 $\pm$ 0.084 & $-$0.0132 & $-$0.0160 to $-$0.0105 & no & yes & yes \\
11 & 0.521 $\pm$ 0.084 & $-$0.0202 & $-$0.0245 to $-$0.0160 & no & no & yes \\
10 & 0.515 $\pm$ 0.083 & $-$0.0258 & $-$0.0311 to $-$0.0205 & no & no & no \\
9 & 0.510 $\pm$ 0.086 & $-$0.0316 & $-$0.0385 to $-$0.0245 & no & no & no \\
8 & 0.505 $\pm$ 0.085 & $-$0.0366 & $-$0.0440 to $-$0.0289 & no & no & no \\
7 & 0.500 $\pm$ 0.085 & $-$0.0412 & $-$0.0488 to $-$0.0331 & no & no & no \\
6 & 0.480 $\pm$ 0.083 & $-$0.0609 & $-$0.0711 to $-$0.0507 & no & no & no \\
5 & 0.428 $\pm$ 0.076 & $-$0.1132 & $-$0.1276 to $-$0.0993 & no & no & no \\
4 & 0.397 $\pm$ 0.071 & $-$0.1444 & $-$0.1590 to $-$0.1306 & no & no & no \\
3 & 0.340 $\pm$ 0.057 & $-$0.2014 & $-$0.2211 to $-$0.1820 & no & no & no \\
2 & 0.299 $\pm$ 0.057 & $-$0.2420 & $-$0.2603 to $-$0.2230 & no & no & no \\
1 & 0.181 $\pm$ 0.040 & $-$0.3602 & $-$0.3831 to $-$0.3368 & no & no & no \\
\bottomrule
\end{tabular}
\end{table*}


\begin{figure*}[t]
\centering
\includegraphics[width=\textwidth]{figures/fig4_sparse_curve.pdf}
\caption{Sparse sensor curve (43 participants). (a) Mean Macro-F1 of models retrained on the top-$k$ electrodes of each fold's cross-fitted ranking, for the four values of $\lambda$ (curves overlap). Band = 95\% cluster-bootstrap CI for $\lambda = 0$. (b) Paired difference to the $k = 16$ model with 95\% CI. Dashed lines mark the prespecified margin ($\delta = 0.02$) and the two sensitivity margins.}
\label{fig:curve}
\end{figure*}

\subsection{H2: attribution-ranked subsets beat random subsets; other rankings match them}
Under the shared comparison model (Table~\ref{tab:baselines}, Fig.~\ref{fig:baselines}), SR-CFA subsets outperformed the mean of random subsets at every tested $k$ (Holm-corrected, 8 tests). The advantage was largest at intermediate sizes: $\HTwoKSixDiff$ Macro-F1 at $k = 6$ (95\% CI $\HTwoKSixLCB$ to $\HTwoKSixUCB$, $p_{\mathrm{Holm}} \HTwoKSixP$, $r_{rb} = \HTwoKSixR$) and $\HTwoKEightDiff$ at $k = 8$. It was small at $k = 4$ ($\HTwoKFourDiff$) and at $k = 12$ ($\HTwoKTwelveDiff$, $p_{\mathrm{Holm}} \HTwoKTwelveP$). Within each participant's own random subsets, the SR-CFA subset had a median percentile of \PctKSix\ at $k = 6$, \PctKEight\ at $k = 8$, \PctKFour\ at $k = 4$ and \PctKTwelve\ at $k = 12$. It reached the 95th percentile for only \PctKAboveNinetyFiveSix\ of 43 participants at $k = 6$ and \PctKAboveNinetyFiveTwelve\ at $k = 12$. Attribution therefore reliably beats the average random subset, but it rarely finds the best subset available for a given person.

In exploratory contrasts, attribution was not better than simpler alternatives. It did not differ from permutation importance at $k = 4$, 8 or 12 (e.g.\ $k = 8$: $\VSPermKEightDiff$, $p \VSPermKEightP$), and it exceeded it only at $k = 6$ ($\VSPermKSixDiff$, $p \VSPermKSixP$). The geometric baseline, which spreads electrodes evenly around both rings without using any data, matched SR-CFA at $k = 4$ ($\VSGeomKFourDiff$, $p \VSGeomKFourP$), $k = 8$ ($\VSGeomKEightDiff$, $p \VSGeomKEightP$) and $k = 12$ ($\VSGeomKTwelveDiff$, $p \VSGeomKTwelveP$). It was worse only at $k = 6$ ($\VSGeomKSixDiff$, $p \VSGeomKSixP$). Mutual information was the weakest ranking at every $k \le 8$.

\begin{table*}[t]
\centering
\small
\caption{Channel-selection methods at matched $k$ (43 participants). All methods use one comparison model (XGBoost, 50 trees, 150k-window outer-training subsample), so values are comparable within the table but lower than Table~\ref{tab:sparse_ni}. H2 $p$ values Holm-corrected over 8 tests; $^\dagger$exploratory, uncorrected.}
\label{tab:baselines}
\begin{tabular}{lcccc}
\toprule
Method & $k$ = 4 & $k$ = 6 & $k$ = 8 & $k$ = 12 \\
\midrule
SR-CFA ($\lambda$=0.5) & 0.383 $\pm$ 0.071 & 0.453 $\pm$ 0.080 & 0.469 $\pm$ 0.081 & 0.485 $\pm$ 0.079 \\
Cross-fitted SHAP ($\lambda$=0) & 0.382 $\pm$ 0.071 & 0.453 $\pm$ 0.080 & 0.468 $\pm$ 0.081 & 0.485 $\pm$ 0.079 \\
Permutation importance & 0.385 $\pm$ 0.070 & 0.440 $\pm$ 0.080 & 0.466 $\pm$ 0.080 & 0.485 $\pm$ 0.079 \\
Geometric farthest-point & 0.383 $\pm$ 0.069 & 0.419 $\pm$ 0.068 & 0.465 $\pm$ 0.072 & 0.478 $\pm$ 0.076 \\
Mutual information & 0.280 $\pm$ 0.049 & 0.360 $\pm$ 0.066 & 0.396 $\pm$ 0.068 & 0.480 $\pm$ 0.079 \\
Random subsets (mean) & 0.366 $\pm$ 0.059 & 0.416 $\pm$ 0.068 & 0.445 $\pm$ 0.071 & 0.478 $\pm$ 0.077 \\
\midrule SR-CFA $-$ random (H2) & +0.0172, $p$=0.007 & +0.0370, $p$=$<$0.001 & +0.0243, $p$=$<$0.001 & +0.0074, $p$=$<$0.001 \\
SR-CFA $-$ perm$^\dagger$ & $-$0.0018, $p$=0.678 & +0.0126, $p$=0.009 & +0.0033, $p$=0.069 & +0.0001, $p$=0.995 \\
SR-CFA $-$ geom$^\dagger$ & $-$0.0005, $p$=0.986 & +0.0341, $p$=$<$0.001 & +0.0047, $p$=0.267 & +0.0070, $p$=0.058 \\
SR-CFA percentile in random draws (median) & 65 (100 draws) & 89 (100 draws) & 82 (100 draws) & 66 (500 draws) \\
\bottomrule
\end{tabular}
\end{table*}


\begin{figure*}[t]
\centering
\includegraphics[width=\textwidth]{figures/fig5_selection_baselines.pdf}
\caption{Selection methods at matched $k$ (43 participants; comparison model shared by all methods). (a) Mean Macro-F1 with 95\% cluster-bootstrap CI. (b) Percentile of the SR-CFA subset within each participant's own random subsets (one point per participant; bar = median; 100 draws, 500 at $k = 12$).}
\label{fig:baselines}
\end{figure*}

\subsection{H3: the stability penalty had no measurable effect}
The instability term was small relative to attribution magnitude: the median $U_c/A_c$ was \UOverAMedian\% (maximum \UOverAMax\%). Penalising it therefore rarely changed a ranking. Against $\lambda = 0$, the top-12 set changed for \TopTwelveDiffersLambdaPtTwoFive, \TopTwelveDiffersLambdaPtFive\ and \TopTwelveDiffersLambdaOne\ of 43 folds at $\lambda = 0.25$, 0.5 and 1. Across all four $k$ tested, at least \HThreeMinTied\ of 43 participants received the identical subset under both rankings. None of the \HThreeNTests\ contrasts was significant after Holm correction (largest mean difference \HThreeMaxAbsDiff; smallest $p_{\mathrm{Holm}} \HThreeMinP$; Supplementary Table~\ref{tab:lambda}). The four curves in Fig.~\ref{fig:curve}a are indistinguishable (maximum spread between $\lambda$ values \LambdaCurveMaxSpread), and $k^*$ was identical under every $\lambda$ and margin. H3 is not supported. In these data, the cross-fitted magnitude ranking is already stable enough that a stability penalty has nothing to correct.

\subsection{H4: the attribution is faithful}
When all three deletion arms were trained under one protocol, the expected ordering held at every retained size (Table~\ref{tab:faithfulness}, Fig.~\ref{fig:faithfulness}). Deleting the highest-ranked electrodes hurt more than deleting random ones, and random deletion hurt more than deleting the lowest-ranked. \FaithNSig\ of the \FaithNTests\ Holm-corrected contrasts were significant. At six retained electrodes, top-removed scored \FaithTopKSix, random \FaithRandomKSix\ and bottom-removed \FaithBottomKSix\ (top $-$ random $\FaithTRKSixDiff$, $p_{\mathrm{Holm}} \FaithTRKSixP$; random $-$ bottom $\FaithRBKSixDiff$, $p_{\mathrm{Holm}} \FaithRBKSixP$). The only non-significant contrast was top- versus bottom-removal at 14 retained electrodes ($\FaithTBKFourteenDiff$, $p_{\mathrm{Holm}} \FaithTBKFourteenP$), where only two electrodes are removed. The gap between arms was largest at intermediate sizes and small at the extremes. With few electrodes left all arms perform poorly; with many left, the array is redundant enough that any two electrodes can be spared.

\begin{figure}[t]
\centering
\includegraphics[width=\columnwidth]{figures/fig7_faithfulness.pdf}
\caption{Deletion-with-retraining test of attribution faithfulness (43 participants; mean and 95\% cluster-bootstrap CI). All arms use the same comparison model and training subsample. Random removal is available where random subsets were sampled ($n$ = 4, 6, 8, 12).}
\label{fig:faithfulness}
\end{figure}

\begin{table*}[t]
\centering
\small
\caption{Deletion (remove-and-retrain) faithfulness test (H4; 43 participants). Every arm uses the same comparison model and training subsample; random removal averages 100 (500 at $k$ = 12) draws per participant and exists only where those draws were sampled. Differences are paired means with Holm-corrected $p$ in parentheses.}
\label{tab:faithfulness}
\begin{tabular}{rcccccc}
\toprule
Kept & Top removed & Random & Bottom removed & Top$-$Bottom & Top$-$Random & Random$-$Bottom \\
\midrule
2 & 0.263 & -- & 0.292 & $-$0.0297 ($<$0.001) & -- & -- \\
4 & 0.351 & 0.366 & 0.382 & $-$0.0314 ($<$0.001) & $-$0.0151 (0.010) & $-$0.0163 (0.013) \\
6 & 0.376 & 0.416 & 0.454 & $-$0.0775 ($<$0.001) & $-$0.0398 ($<$0.001) & $-$0.0377 ($<$0.001) \\
8 & 0.408 & 0.445 & 0.469 & $-$0.0618 ($<$0.001) & $-$0.0377 ($<$0.001) & $-$0.0242 ($<$0.001) \\
10 & 0.445 & -- & 0.475 & $-$0.0297 ($<$0.001) & -- & -- \\
12 & 0.455 & 0.478 & 0.485 & $-$0.0302 ($<$0.001) & $-$0.0228 ($<$0.001) & $-$0.0074 ($<$0.001) \\
14 & 0.487 & -- & 0.491 & $-$0.0038 (0.150) & -- & -- \\
\bottomrule
\end{tabular}
\end{table*}


\subsection{Stability of the ranking across people}
Rankings computed from individual participants' attributions agreed closely with each other (Table~\ref{tab:stability}): mean Spearman $\rho$ = \StabSubjRho\ $\pm$ \StabSubjRhoSD\ over \StabSubjPairs\ pairs, and Jaccard overlap of the top six = \StabSubjJacSix. Rankings from disjoint inner folds agreed at $\rho$ = \StabInnerRho. The outer-fold agreement ($\rho$ = \StabOuterRho) is an upper bound only, because those training sets share about 98\% of their participants. Across folds, the highest- and lowest-ranked electrodes kept their positions, and the rank ranges were widest in the middle of the ordering (Fig.~\ref{fig:attribution}b).

\begin{table*}[t]
\centering
\small
\caption{Agreement of channel rankings ($\lambda = 0$). Values are means over pairs.}
\label{tab:stability}
\begin{tabular}{lcccccc}
\toprule
Comparison & Pairs & Spearman $\rho$ & Kendall $\tau$ & Jaccard@4 & Jaccard@6 & Jaccard@8 \\
\midrule
Per-subject attribution (between people) & 903 & 0.927 $\pm$ 0.038 & 0.801 & 0.752 & 0.905 & 0.803 \\
Inner folds (disjoint subject groups) & 430 & 0.948 $\pm$ 0.021 & 0.840 & 0.816 & 0.992 & 0.838 \\
Outer folds (upper bound$^*$) & 903 & 0.987 $\pm$ 0.008 & 0.943 & 0.888 & 1.000 & 0.923 \\
\bottomrule
\end{tabular}
\par\smallskip\parbox{0.95\linewidth}{\footnotesize $^*$Outer training sets share about 98\% of their participants, so their agreement is an upper bound, not an estimate of between-subject stability.}
\end{table*}


\begin{figure*}[t]
\centering
\includegraphics[width=\textwidth]{figures/fig6_attribution.pdf}
\caption{Cross-fitted channel attribution ($\lambda = 0$, 43 outer folds). (a) Mean share of attribution per electrode, drawn on the $2 \times 8$ montage (ring~2 electrodes are aligned with, and 2~cm from, ring~1 electrodes at the same position). (b) Rank of each electrode across folds (median and range).}
\label{fig:attribution}
\end{figure*}

\subsection{Exploratory analyses}

\begin{figure*}[t]
\centering
\includegraphics[width=\textwidth]{figures/fig8_generalisation.pdf}
\caption{Exploratory analyses. (a) Cross-session: each point is one held-out participant (bar = median); all conditions are also subject-independent. (b) NinaPro DB5 sparse curve (10 subjects; mean and 95\% cluster-bootstrap CI; dotted line = chance). (c) GRABMyo bipolar-pair sparse curve (43 participants; $\lambda = 0$).}
\label{fig:generalisation}
\end{figure*}

\paragraph{Cross-session}
Training on session~1 and testing on sessions 1, 2 and 3 gave \CSOneSOne, \CSOneSTwo\ and \CSOneSThree\ (Table~\ref{tab:cross_session}, Fig.~\ref{fig:generalisation}a). Performance on later sessions was not lower than on the training session. The session-2 score was higher ($\CSSTwoVsSOneDiff$, 95\% CI $\CSSTwoVsSOneLCB$ to $\CSSTwoVsSOneUCB$, uncorrected $p \CSSTwoVsSOneP$), and session 3 did not differ ($\CSSThreeVsSOneDiff$, $p \CSSThreeVsSOneP$). Because every condition is also subject-independent, this does not show that sessions are interchangeable. It shows that, once the cross-user gap is present, no additional cross-day loss was detectable. Training on sessions 1 and 2 improved session-3 performance over training on session 1 alone ($\CSTwoVsOneDiff$, 95\% CI $\CSTwoVsOneLCB$ to $\CSTwoVsOneUCB$, $p \CSTwoVsOneP$). This design cannot separate day diversity from the doubled amount of training data.

\paragraph{Bipolar montage}
With the eight derived bipolar pairs as the selection unit, full-array LOSO Macro-F1 was \BipFull\ $\pm$ \BipFullSD, lower than the monopolar array (Fig.~\ref{fig:generalisation}c, Supplementary Table~\ref{tab:replication}). No pair could be removed within the primary margin. With seven pairs, $\bar\Delta_7 = \BipDeltaKSeven$, and its lower bound ($\BipDeltaKSevenLCB$) fell just outside $-0.02$. It did so under every one of 200 alternative bootstrap seeds (range $\BipSeedLCBMinSeven$ to $\BipSeedLCBMaxSeven$). $k^*$ was therefore \BipKStarMarginPtTwo\ at $\delta = 0.02$, \BipKStarMarginPtOne\ at 0.01, and \BipKStarMarginPtThree\ pairs at 0.03, identically for all $\lambda$. The ranking was again stable: pairs \BipTopThreeSet\ formed the top three in \BipTopThreeFolds\ of 43 folds. In deployable sensor units, then, the array showed no redundancy within the prespecified tolerance.

\paragraph{NinaPro DB5}
On DB5 the full 16-channel LOSO Macro-F1 was \DBFiveFull\ $\pm$ \DBFiveFullSD\ (12 movements; chance about 0.083), far below GRABMyo and with wide between-subject spread (\DBFiveFullMin--\DBFiveFullMax) (Fig.~\ref{fig:generalisation}b, Supplementary Table~\ref{tab:replication}). Under the frozen rule, $k^* = \DBFiveKStarMarginPtTwo$ at $\delta = 0.02$ ($\bar\Delta_{12} = \DBFiveDeltaKTwelve$, lower bound $\DBFiveDeltaKTwelveLCB$). This result is fragile. The lower bound lies on the margin, and $k = 12$ was non-inferior under \DBFiveSeedFrac\% of \DBFiveSeedN\ alternative bootstrap seeds (lower bound range $\DBFiveSeedLCBMin$ to $\DBFiveSeedLCBMax$). Moreover, $k = 13$, 14 and 15 all failed (lower bounds $\DBFiveDeltaKThirteenLCB$, $\DBFiveDeltaKFourteenLCB$, $\DBFiveDeltaKFifteenLCB$). A fixed-sequence procedure starting at $k = 15$ would stop immediately and report no non-inferior subset. With ten subjects, DB5 does not support a claim that a quarter of the channels can be removed. The procedure itself transferred without modification, and the ranked curve had the same monotone shape.

\paragraph{Fixed layout (post hoc)}
One 12-electrode layout (removing electrodes \FixedDropped) was derived on \FixedDev\ participants and applied unchanged to \FixedTest\ others. It scored $\FixedDiff$ against the 16-electrode model trained on the same participants (95\% CI $\FixedLCB$ to $\FixedUCB$). This lies within the margin, but it is one post-hoc split and is reported only as a consistency check on the LOSO result.

\paragraph{Feature redundancy}
Removing IEMG, an exact multiple of MAV, changed Macro-F1 by $\RedundDiff$ (95\% CI $\RedundLCB$ to $\RedundUCB$; 43 folds, comparison model). Channel rankings were almost unchanged: the median per-fold Spearman $\rho$ between rankings with and without IEMG was \RedundRho, and the top-12 set was identical in \RedundTopTwelveSame\ of 43 folds. The duplicated feature therefore did not distort the channel ranking.
